\section{Results Across 175 SPARC Galaxies}

% TODO: Port text/equations into each subsection; verify against the PDF.

\subsection{Global Fit Performance}
The boundary-layer response model with a single amplitude parameter \emph{(see PDF: $Q$)} provides substantial improvement over the baryons-only prediction across the SPARC sample. Using a \emph{(see PDF: $\chi^2$-based improvement metric)}, we find that the response-sector closure delivers strong improvements for the large majority of galaxies while remaining parsimonious (one fitted amplitude plus a data-derived transition radius).

\emph{(Table from PDF: global fit performance summary across 175 galaxies.)}

\subsection{Representative Galaxy Diagnostics}
To illustrate performance across the galaxy mass spectrum, we report representative six-panel diagnostics spanning gas-dominated dwarfs to massive spirals. The figures will be inserted during the PDF-faithful port (see PDF for current canonical visuals).

Representative cases discussed in the source draft include:
\begin{itemize}
\item NGC 2403 (intermediate-mass spiral);
\item DDO 154 (gas-dominated dwarf irregular);
\item NGC 2841 (massive high-surface-brightness spiral);
\item UGC 06614 (low-surface-brightness system).
\end{itemize}

\emph{(Figures from PDF: six-panel diagnostics for these galaxies.)}

\subsection{Emergent Baryonic Tully–Fisher Relation}
A critical test of any dark-matter alternative is whether it recovers the baryonic Tully--Fisher relation (BTFR): the tight empirical correlation between baryonic mass and asymptotic rotation velocity \citep{McGaugh2000}. In our framework, the BTFR is not imposed as a prior but emerges as a consequence of the boundary-layer closure.

\emph{(See PDF: correlation statement between fitted amplitude and BTFR proxy, with Spearman $\rho$ and $p$-value.)}

\subsection{Fitted vs. Robust Amplitude Comparison}
A key robustness check compares the model-fitted amplitude \emph{(see PDF: $Q_{\mathrm{fit}}$)} (from global minimization) with the independently computed \emph{(see PDF: $Q_{\mathrm{est}}$)} (from Huber robust estimation on the outer velocity deficit). Across 175 galaxies, the rank correlation is strong, supporting the claim that the fitted amplitude reflects outer-region physics rather than overfitting to inner structure.

\emph{(See PDF: Spearman $\rho$, fraction with $Q_{\mathrm{est}}/Q_{\mathrm{fit}} < 1$, and median ratio.)}

\subsection{Mass-to-Light Ratio Sensitivity}
Sweeping $\Upsilon_\star$ across \emph{(see PDF: the tested grid)} demonstrates that response-sector inferences are robust to plausible photometric systematics. While individual galaxy amplitudes shift with $\Upsilon_\star$, the global performance metrics and BTFR scaling remain stable.

\emph{(See PDF: median fractional change in amplitude across the $\Upsilon_\star$ sweep and the dynamic range of $Q$.)}

\subsection{Oscillatory Signatures and Negative-Q Galaxies}
Two galaxies \emph{(see PDF: named cases)} exhibit negative \emph{(see PDF: $Q$)} values---cases where the observed velocity falls below the baryonic prediction in the outer regions. Within the response-sector framework, these are interpreted as rarefaction-phase sampling of an oscillatory boundary-layer response: the response includes both compression ($\Delta a > 0$) and rarefaction ($\Delta a < 0$) phases, and some galaxies' observed radial coverage samples primarily the rarefaction phase.

More broadly, \emph{(see PDF: count and percentage)} of galaxies show at least one sign change in their velocity-deficit profiles. The characteristic oscillation wavelength scales with galaxy size (median \emph{(see PDF)}), consistent with standing-wave configurations in the scalar sector of the metric field that are distinct from transverse spin-2 gravitational waves.
